\documentclass[11pt,a4paper]{article}
\usepackage[utf8]{inputenc}
\usepackage[T2A]{fontenc}
\usepackage[russian,english]{babel}
\usepackage[margin=22mm]{geometry}
\usepackage{booktabs,amsmath,hyperref}
\hypersetup{colorlinks=true,urlcolor=blue}
\title{Пилот Al$_{32}$: план реального DFT-датасета\newline для DPA4 и GRACE}
\author{MLIP-nucleation project}
\date{18 октября 2026 г.}
\begin{document}
\selectlanguage{russian}
\maketitle

\section{Решение и границы применимости}
Выбран однокомпонентный Al в периодических ячейках из 32 атомов.
Это ограниченный пилот адаптации foundation models к собственному DFT,
а не окончательный выбор сплава и инокулятора. Гипотеза о причинной роли
межфазного порядка остаётся непроверенной. Интерес к Al связан с исследованиями
преднуклеационных слоёв жидкого Al на подложках~\cite{layering} и
композиционного шаблонирования интерметаллидов~\cite{templating}.
Новизна выбранной системы не заявляется.

\begin{center}
\begin{tabular}{p{22mm}p{125mm}}
\toprule
Кандидат & Решение для первого пилота\\
\midrule
Al & Один элемент, немагнитный расчёт, связь с исходным исследованием;
выбран для проверки всей цепочки меток.\\
Cu & Допустимая простая альтернатива; менее прямая связь с исходным brief.\\
Ni & Отложен: требуется согласовать описание магнитного состояния.\\
Al--Ti/TiB$_2$ & Научно ближе к сегрегации, но нужны несколько химических
окружений, терминации и более крупные интерфейсные ячейки.\\
\bottomrule
\end{tabular}
\end{center}
Оба загруженных checkpoint содержат тип Al. Наличие элемента в type map
не доказывает точность. Полный набор для исследования должен отдельно покрыть
поверхности, интерфейсы, сегрегацию, дефекты и достаточно большие жидкие ячейки.
Этот пилот не определяет температуру плавления, скорость зарождения или барьер.

\section{Конфигурации и заранее заданное разбиение}
Запланировано 40 конфигураций: 9 FCC со сжатием/расширением, сдвигом и
смещениями атомов; 24 конфигурации из восьми независимо инициализированных
траекторий EMT; 7 HCP при разных объёмах. Каждая ячейка содержит 32 атома.
Все производные одного FCC-прототипа находятся только в train, HCP ---
только в test как проверка переноса на другой тип упаковки.
Из восьми траекторий четыре целиком идут в train, две в validation, две в test.
Итого train/validation/test: 21/6/13; в test отдельно анализируются 6
разупорядоченных и 7 HCP конфигураций. Соседние кадры не являются независимыми
наблюдениями; статистику следует группировать по родительским траекториям.
Все траектории начинаются с общего FCC-прототипа; независимые seed не означают
независимые химические системы или отсутствие сходных структур в pretraining.

EMT применяется исключительно для предложения геометрий~\cite{emt}.
План: 1000 шагов при 2500 K, 1000 при целевой температуре 1000 или 1400 K,
затем три кадра с интервалом 100 шагов; шаг 2 fs, Langevin NVT.
Объёмы 17.5 и 18.5~\AA$^3$/атом. Это не AIMD и не доказательство
равновесной жидкости DFT. Сохраняются seed, шаг, температура, минимальное
расстояние и диагностическая амплитуда FCC-пика. Состояние разупорядочения
не следует отождествлять с термодинамической фазой по одной метрике.
Метки EMT не экспортируются.

\section{Протокол VASP и ворота сходимости}
VASP 6.6.1 CPU/MPI, PBE, PAW Al из пользовательской библиотеки 6.4;
точный SHA-256 POTCAR закреплён в конфигурации и каждом расчёте.
Статические расчёты: \texttt{ISPIN=1}, \texttt{ISYM=0},
\texttt{PREC=Accurate}, \texttt{LREAL=.FALSE.}, \texttt{LASPH=.TRUE.},
\texttt{ADDGRID=.TRUE.}, \texttt{ALGO=Normal},
\texttt{EDIFF=1E-7}, \texttt{NELM=160}, \texttt{NSW=0}.

На двух заранее выбранных тренировочных конфигурациях проверяются
\texttt{ENCUT}=350/450/550 eV, Gamma-сетки $3^3/4^3/5^3$,
\texttt{ISMEAR=0}, $\sigma=0.05/0.10/0.20$ eV; всего 16 single points.
Опорная точка: 550 eV, $5^3$, 0.05 eV; кандидат для серии:
450 eV, $4^3$, 0.05 eV.
Кандидат, отдельные 450 eV и $4^3$, а также изменение $\sigma$ до 0.10 eV
должны отличаться от опоры не более чем на 2 meV/атом,
0.02 eV/\AA\ по любой компоненте силы и
0.001 eV/\AA$^3$ по любой компоненте напряжения.
Пределы заданы до результатов. Сходимость требуется проверять для интересующей
величины, а не только полной энергии~\cite{encut}.
При неудаче серия блокируется; повышение опорной точности или бюджета требует
нового плана. Проверка двух структур является локальной оценкой ошибки,
не доказательством сходимости всей будущей химической области.

\section{Контракт меток и проверка экспорта}
Энергия обучения --- свободная энергия \texttt{TOTEN} в eV на ячейку.
Энергия $\sigma\to0$ сохраняется отдельно. При конечном размытии силы и
напряжения согласованы со свободной энергией~\cite{smearing}.
Силы: eV/\AA. Напряжение: растяжение положительно,
eV/\AA$^3$, Voigt $xx,yy,zz,yz,xz,xy$.
VASP pressure-положительный тензор в kbar меняет знак при переходе к ASE;
конверсия независимо сверяется с XML. Virial $W=-V\sigma$ в eV,
полная матрица $3\times3$, row-major.

Парсер требует полный XML, ровно один статический шаг, достижение EDIFF до
NELM, нормальное завершение VASP, конечные числа и правильные размерности.
Неуспешный расчёт не становится меткой. Экспорты: ASE/GRACE extxyz и
DeepMD NPY с type map Al; идентификаторы кадров и SHA-256 сохраняются.
Проверяется чтение extxyz обратно и согласованность знака virial.

\section{Совместимость foundation models и будущая проверка}
OMat24 использует PBE/PBE+U и PAW версии 5.4; авторы предупреждают о
несовместимости абсолютных энергий при смешении настроек~\cite{omat}.
Для металлического Al Hubbard U не нужен. Совпадение конкретных PAW-файлов
5.4 и 6.4 побайтно не установлено; считать их эквивалентными нельзя.
Все новые метки имеют один локальный протокол; OMat24 replay пока не добавляется.
Исходные DPA4-Mini-OMat24 и GRACE-1L-OMAT веса закреплены в
\texttt{models/foundation/}; синтетические дообученные веса не используются.

Для будущего обучения любой элементный сдвиг энергии определяется только
на train. Отчёт должен показывать как исходные, так и сдвинутые ошибки:
сдвиг не устраняет ошибки сил/напряжений и не доказывает эквивалентность PAW.
Новые seed не исключают сходства с FCC/HCP и другими структурами OMat24;
этот набор не является доказанно независимым от pretraining benchmark.
Validation применяется для выбора checkpoint, test --- только для итоговой
оценки, без подбора sampling, гиперпараметров или константы энергии.
Предварительные ориентиры пилота: RMSE энергии 10 meV/атом после train-сдвига,
сил 0.10 eV/\AA, напряжения 0.01 eV/\AA$^3$; итог приводится отдельно по
разупорядоченным структурам и HCP. Эти ориентиры не квалифицируют модель для
расчёта барьера зарождения. Дополнительный контроль --- форма HCP $E(V)$
и разность энергии HCP относительно FCC при сопоставимых объёмах.

\section{Выполнение и сохранность}
План ресурсов: \texttt{apollo}, account \texttt{imet-2026-2}, QoS normal,
один узел, 16 MPI ranks, 32 GiB RAM; максимум 2 h калибровки и 6 h разметки,
то есть 128 CPU$\cdot$h зарезервированного времени. GPU не используется.
Подача разрешена только после явного подтверждения этого бюджета пользователем.
Перед подачей сохраняются run ID, хеши и commit; job ID записывается сразу.
Неопределённый результат подачи требует сверки scheduler, а не повтора.

Публичные входы, настройки, разбиение и небольшие DFT-метки идут в GitHub.
POTCAR и сырые VASP-выходы сохраняются в частном архиве в основном home кластера
и фактической проверенной копии на рабочем компьютере. SHA-256 проверяется
как для архива, так и для его членов. Срок хранения и институциональное
резервирование пока не подтверждены; scratch не считается архивом.

\paragraph{Статус.} Подготовлены конфигурация и исполняемый workflow.
Сходимость и наличие принятых DFT-меток определяются только записями run,
а не этим планом. До завершения расчётов число принятых меток равно нулю.

\begin{thebibliography}{9}
\bibitem{layering}\href{https://www.nature.com/articles/s42004-018-0104-1}{Atomistics of pre-nucleation layering of liquid metals at the interface with poor nucleants}, Communications Chemistry (2018).
\bibitem{templating}Z. Que et al., \href{https://www.nature.com/articles/s41598-024-59709-w}{Composition templating for heterogeneous nucleation of intermetallic compounds}, Scientific Reports 14, 8968 (2024).
\bibitem{emt}\href{https://docs.ase-lib.org/_modules/ase/calculators/emt.html}{ASE EMT implementation and parameters} (accessed 2026-10-18).
\bibitem{encut}\href{https://vasp.at/wiki/ENCUT}{VASP Wiki: ENCUT} (accessed 2026-10-18).
\bibitem{smearing}\href{https://vasp.at/wiki/index.php/Smearing_technique}{VASP Wiki: Smearing technique} (accessed 2026-10-18).
\bibitem{omat}L. Barroso-Luque et al., \href{https://arxiv.org/html/2410.12771v1}{Open Materials 2024 (OMat24) Inorganic Materials Dataset and Models}, arXiv:2410.12771, Sec. 2.2.2 (2024).
\end{thebibliography}
\end{document}
